% !TeX root = ../informe.tex
% ===========================================================================
%  Resumen ejecutivo
%  Extensión sugerida: 150-250 palabras, en un solo párrafo.
%  Debe responder: qué se construyó, con qué arquitectura, y qué resultado
%  se obtuvo. Se escribe al final, cuando el resto del informe ya está listo.
% ===========================================================================

\section*{Resumen}
\addcontentsline{toc}{section}{Resumen}

Se desarrolló un sistema web para la reserva de canchas deportivas de tenis,
básquet y pádel, que permite consultar disponibilidad, crear y cancelar reservas,
administrar canchas y usuarios, y visualizar reportes básicos. La aplicación se
organizó mediante una arquitectura basada en microfrontends y microservicios. En
el frontend se implementó un shell junto con tres microfrontends remotos
utilizando React y Module Federation, mientras que el backend se dividió en los
microservicios ms-usuarios, ms-canchas, ms-reservas y ms-reportes desarrollados
con Spring Boot. La persistencia se organizó en PostgreSQL mediante bases de
datos independientes para los microservicios que requieren almacenar
información, evitando el acceso directo a los datos administrados por otros
servicios. También se implementaron las reglas relacionadas con permisos,
cancelaciones, estados de reserva, límite de reservas activas y control de
solapamiento de horarios. Para este último caso, la validación realizada por
ms-reservas se complementó con una restricción de exclusión en PostgreSQL para
evitar conflictos ante solicitudes concurrentes. Los componentes del sistema se
ejecutan de forma conjunta mediante Docker Compose. Como resultado, se
implementaron las funcionalidades definidas para el usuario final y el
administrador, y los seis criterios de aceptación establecidos para el proyecto
se encuentran cubiertos por la implementación y las pruebas realizadas.

\vspace{0.5cm}
\noindent\textbf{Palabras clave:} microservicios, microfrontends,
Module Federation, Spring Boot, PostgreSQL, Docker.